\chapter{The investment case and its alternatives}
\label{r3:investment}

A common causal product earns its cost by improving repeated scientific
work. The appropriate alternative is a capable quantitative team composing
the same open-source methods with ordinary engineering, documentation and
review. Such a team can produce good causal studies without a new platform.
The investment question is whether the shared representations and qualified
operations reduce the cost of equivalent work or enable additional valuable
decisions, after construction, support and adoption are included.

\section{Define equivalent output before pricing reuse}
\label{r3:equivalent-output}

An equivalent completed study has a reconstructed population and observation
clock, a clear intervention and financial outcome, a defensible identification
argument or useful nonidentification conclusion, appropriate inference,
reproducible artifacts and review. The comparison must include negative
results. An early finding that a requested target is unidentified can save
substantial work; forcing every study to return a numerical effect would
reward scientific error.

The proposed service would help analysts reconstruct what was known when
a decision was made, keep the population and outcome consistent, document
why an effect is identified, and carry uncertainty into the financial
comparison. It would reuse existing numerical methods where their
assumptions fit. Each business would still specify its products, costs and
constraints. A mortgage policy, an individual collections decision and a
macroeconomic shock need not share the same economic quantity of interest.

Comparator costs should include the same required outputs and review
quality. Reduced effort from omitting uncertainty or failing to investigate
weak support is not a saving. Conversely, building a generic model interface
for a one-off study can cost more than a small, clearly documented expert
analysis. The portfolio of genuinely repeated requests, not the total
number of database rows, determines the economic opportunity.

\section{An auditable cost model}
\label{r3:cost-model}

For alternative $j\in\{E,P\}$, expert composition or shared product,
let $K_j$ be initial construction cost, $F_j$ annual shared cost,
$c_j$ marginal cost per equivalent completed study, and $n_t$ completed
studies in year $t$. With horizon $H$ and discount rate $r$,
\begin{equation}
 C_j=K_j+\sum_{t=1}^{H}\frac{F_j+c_jn_t}{(1+r)^t}.
 \label{r3:present-cost}
\end{equation}
The categories are disjoint: method maintenance and common infrastructure
belong in $F_j$, while study-specific preparation, fitting, review and
rework belong in $c_j$. Failed-study costs must be included in one
category; a completed-study denominator cannot make them disappear.

Use the following illustrative USD-million assumptions to expose the
break-even logic:
\[
 (K_E,F_E,c_E)=(1,.6,.15),\qquad
 (K_P,F_P,c_P)=(4,1.8,.06),\qquad H=3,\quad r=.10.
\]
These are planning inputs, not observed bank prices, measured productivity
or enterprise vendor quotations. Let
$a_H=\sum_{t=1}^H(1+r)^{-t}$. At constant annual volume $n$, the cost-only
threshold, when $c_E>c_P$, is
\begin{equation}
 n^*=\frac{(K_P-K_E)/a_H+(F_P-F_E)}{c_E-c_P}.
 \label{r3:reuse-threshold}
\end{equation}
If the marginal saving is zero or negative, this particular reuse argument
does not establish a positive-volume crossover. Under the stipulated
numbers, $a_H=2.48685$ and $n^*=26.74$ studies annually.

\begin{table}[htbp]\centering\small
\begin{tabular}{@{}lrrr@{}}
\toprule Scenario & Expert cost & Platform cost & Platform saving\\\midrule
Low reuse & 5.476 & 9.670 & -4.194\\
Moderate reuse & 11.445 & 12.057 & -0.613\\
High reuse & 15.921 & 13.848 & 2.073\\
Adoption ramp & 8.402 & 10.840 & -2.438\\
\bottomrule\end{tabular}

\caption{Three-year present cost in illustrative USD millions. Annual
volumes are 8, 24 and 36, or the ramp $(4,16,30)$. Both alternatives
deliver the same specified scientific outputs.}
\label{r3:reuse-economics}
\end{table}

At eight annual studies the product costs about USD\,4.19 million more.
At 24 it still costs about USD\,.61 million more. At 36 it saves about
USD\,2.07 million. The ramp leaves a USD\,2.44 million premium even
though final-year volume exceeds the constant-volume threshold:
construction occurs before most reuse. Halving marginal savings doubles
the threshold to 53.47 studies; increasing the annual shared-cost
premium to USD\,2 million raises it to 35.63.

\begin{figure}[htbp]\centering
\includegraphics[width=.88\textwidth]{derivations/reuse_economics.pdf}
\caption{The crossover follows from the declared account. It is not
a forecast of demand or evidence that the shared layer already saves effort.}
\end{figure}

The simple model holds study mix fixed. A more realistic portfolio has
classes $k$ with costs $c_{jk}$ and volumes $n_{tk}$, replacing
$c_jn_t$ by $\sum_k c_{jk}n_{tk}$. A few expensive bespoke latent
models may have little shared work with many routine rollouts.
Reuse should therefore be measured by operation and data contract, not
assumed from a common causal-inference label.

\section{Separate business value from construction savings}
\label{r3:incremental-value}

Let $D_t$ be incremental realized decision value attributable to the
product relative to the equivalent expert workflow. Then
\begin{equation}
 \Delta\operatorname{NPV}
 =C_E-C_P+\sum_{t=1}^{H}\frac{D_t}{(1+r)^t}.
 \label{r3:incremental-npv}
\end{equation}
The first term already counts workflow savings. Released analyst time
must not be added again as business value. If it permits extra studies,
their costs, capacity use, outcome lags and incremental decisions
belong in the expanded account.

At 24 annual studies the illustrative cost premium requires about
USD\,.246 million of additional constant annual decision value.
The adoption ramp requires about USD\,.980 million. The size of a
bank's balance sheet does not establish either amount. The expert
comparator may already identify the same profitable policy or reject
the same loss-making offer.

The deposit study illustrates the attribution problem. Finding a
negative contribution can prevent a loss only if the finding changes
what would otherwise have been done. The prevented loss cannot be
credited to the product merely because its report contains the estimate.
Likewise the simulated collections gain is a property of the declared
environment; it supplies no observed revenue to this investment account.
Latent-state research can be worthwhile even while operational value
remains uncertain.

Costs and gains also mature on different clocks. Data access,
implementation and review are paid before customer outcomes are
observed. A six-month retention measure may mature during an early
evaluation; credit losses or a persistent pricing response may not.
A continuation model supplies a conditional forecast for those outcomes,
not an observation. Report the cash timing and model dependence explicitly.

\section{Measure the common-product hypothesis}
\label{r3:adoption-experiment}

The complete-workflow measurements in Section~\ref{s4:workflow-results}
refine the resource account but do not estimate its labor terms. A response
workflow over 3.85 million versioned records takes about 3.38 seconds on the
recorded host; a 3,200-customer latent workflow with 19 bootstrap refits takes
about 48.27 seconds. Both the product and expert-composition arms can use
these same numerical operations. Pricing one arm from a bare estimator call
and the other from reconstruction plus inference would be an unequal-output
comparison. The illustrative per-study dollar inputs above therefore remain
unchanged pending matched measurements of preparation, review and rework.

The first tranche should compare matched tasks in the intervention,
sequential-decision and dynamic-systems workspaces. Agree their outputs,
scientific scope and reviewer rubric before measuring effort. Give the
expert-composition arm the same sources, data access and expertise, and
allow sensible reusable templates. Counterbalance task order where
possible to reduce learning effects.

Measure analyst and reviewer hours, calendar delay, machine cost,
failed starts, rework and unresolved material findings separately.
Ask a second analyst to reconstruct the population, target and financial
conclusion from the exported artifacts. Deliberately introduced clock,
weight and sign errors can test whether review finds the defects.
This tests the proposed product's specific value proposition, rather
than whether a polished interface is pleasant.

A planning target of 30\% lower median analyst/reviewer effort can be
useful if review quality and output scope are held fixed. It remains
a hypothesis until measured across repeated uses and teams. A successful
first-author demonstration is insufficient evidence of institutional
reuse. Conversely, more early rejections can be a benefit if they
correctly avoid expensive unidentified studies.

Progression can be written as
\[
 n_t^{\rm usable}=n_t^{\rm started}
 p_t^{\rm data}p_t^{\rm design}p_t^{\rm execution}p_t^{\rm review},
\]
where each probability conditions on reaching the preceding stage.
This is the chain rule, not an independence assumption. It separates
volume aspirations from usable output and makes adoption bottlenecks
visible. Charges for unsuccessful studies remain in the cost account.

\section{A staged construction budget and stopping rules}
\label{r3:investment-sequence}

An illustrative first tranche assigns six people for six months:
one product/economics lead, two causal/econometric scientists, two
data/platform engineers and one application engineer. At USD\,25,000
per loaded person-month, staff cost is USD\,900,000; USD\,50,000
compute/tooling and 20\% contingency give USD\,1.14 million.
This tranche is included in the illustrative USD\,4 million initial
product cost, not added to it. Data-owner, independent validation,
security and application-owner participation are unpriced and must
be estimated consistently for both alternatives before a bank funds work.

The tranche should implement the shared observation/request/evidence
objects and qualify a narrow adjusted/cohort and observed-history
policy path. Controlled-state fitting and temporal discovery should
enter the architecture immediately through finite, explicitly conditional
research operations. Six people cannot qualify every method in the
44-project survey for production in six months.

The next tranche depends on concrete outcomes. Continue if repeated
requests share observation and scientific contracts, reviewers can
reconstruct results, the required estimators/inference pass their
acceptance exercises, and measured effort supports the economic case.
Narrow the scope if most work remains application-specific. Retain
expert composition when reuse is low. Stop a particular policy analysis
when its data cannot distinguish materially different decisions and
additional evidence is not worth its cost.

These are different decisions: discontinuing an unidentified analysis,
deferring a numerical adapter, and declining investment in the shared
product should not be collapsed into a single project-failure flag.
The source archive and scientific representations remain useful even
if one deployment case is rejected.

\section{An independent offering must price institutional variation}
\label{r3:commercial-product}

An independent product can distribute the scientific core as an embedded
library and provide the three workspaces with deployment, support and
upgrade services. A design-partner engagement is a plausible first
commercial form: a named study portfolio, paid integration scope,
reproducible evidence export and a recurring support boundary.
No willingness-to-pay, sales-cycle or customer-retention estimate is
established by this proposal.

Each institution adds semantic mapping, legal-entity structure, observation
systems, scientific review and operational requirements. Code reuse
reduces some common work but does not remove those costs. A commercial
comparison should price onboarding, customer-specific adapters, method
maintenance, support and validation participation, and state which
customizations remain owned by the customer.

Procurement should compare products on the same reconstructed studies,
supported identification routes, uncertainty objects, data-locality
requirements and export/replay obligations. Repository activity or an
enterprise demonstration is not evidence that these scientific
operations are available. The source review establishes open-source
construction alternatives; it does not assert uninspected private-vendor
capabilities or prices.

For a banking group, the defensible first investment is therefore a
bounded common scientific layer tied to named repeated decisions.
Its expansion should follow measured reuse and preserved analytical
quality. The mathematical distinctions developed throughout the
monograph determine which work can be shared and what the investment
is actually buying.
